\section{Discussion}
\label{sec:discussion}

The main result is not that one backend wins. The main result is that \emph{checkpoint format} determines which backend wins, and this determination is consistent across model scales.

\subsection{Mechanism: why format changes the answer}

SafeTensors and ServerlessLLM-style partitioning create different I/O problems. SafeTensors exposes large contiguous reads after metadata parsing. For small checkpoints, a practical sync backend has low overhead and enough parallelism; the gap between sync and async is small. As size and shard count grow, async's ability to overlap I/O with scheduling and batching provides increasing returns.

ServerlessLLM-style partitioning changes the request geometry. The loader reconstructs tensors from indexed partition files, so grouping, coalescing, and range batching dominate. Sync's sequential submission model pays a per-request cost that accumulates across the many ranges. Async's concurrent scheduling amortizes this cost. The SafeTensors intuition that ``sync can be competitive for small loads'' does not transfer: sync never wins in the ServerlessLLM matrix. The format, not the API name, explains the reversal.

\subsection{Integration: why rankings move across layers}

Rust and vLLM answer different questions. Rust isolates backend mechanics. vLLM measures serving startup with CUDA context setup, memory pools, runtime configuration, and generation state; its path includes Python integration overhead through Tensora's Python bindings. These costs can dilute a pure I/O advantage or reorder two close backends.

The vLLM results do not reproduce Rust rankings row-for-row because vLLM adds substantial initialization overhead that compresses proportional backend spreads. Yet the regime structure survives: backend-aware loading consistently reduces cold initialization time relative to vLLM's native loader, and the advantage grows with model size. This is not a flaw in the methodology; it is the deployment reality. A systems paper about loading should show both mechanism and propagation.

\subsection{The adaptive default backend}

The adaptive default backend selects sync or async based on checkpoint layout, model scale, and platform capability. On Modal (where \texttt{io\_uring} is unavailable), it primarily selects async for SafeTensors and async for ServerlessLLM. The default backend consistently matches or approaches the best explicit backend at each cell, validating that the regime structure is predictable enough for automatic selection.

\subsection{External baseline context}

The hf-native baseline provides context for the absolute performance of Tensora's backends. The upstream \texttt{safetensors} crate uses a simple \texttt{std::fs::read}-based loading path without backend-aware scheduling. Tensora's async backend consistently outperforms hf-native at mid-to-large scales, demonstrating that explicit backend selection provides practical value over the default upstream path.

\subsection{Implications}

For practitioners, evaluate backend choice relative to checkpoint format and scale rather than adopting a fixed submission policy. When deploying on restricted environments (e.g., gVisor containers where \texttt{io\_uring} is unavailable), async provides consistent advantages over sync across both formats.

For researchers, report loading results with explicit format and backend attribution. A paper that reports only one layout or one backend risks explaining the wrong phenomenon: the same I/O API that wins on contiguous reads may lose on partitioned access.
